\section{Part 2：标准 LLM 并行 primitives}

前面的网络基础确定了通信成本，本部分开始逐一拆解“切什么”。每种 primitive 都会改变参数、梯度、optimizer state、activation、token 或 expert 的归属，并引入相应 collective；阅读时始终用同一张账本比较单卡显存、通信量、global batch 限制、调度空泡与实现复杂度。


\singleslide{slides-images/slide-014.jpg}{Part 2 路线图：data/ZeRO、pipeline/tensor model parallel，以及 sequence activation parallel。}{14}

Slide 14 先按“切分对象”组织术语：data parallel 切 batch，ZeRO 进一步切训练状态；pipeline 按层切 depth，tensor 按矩阵切 width；sequence parallel 切 activation 的 token 轴。后面所有组合方案都可以还原到这些轴，不能只凭缩写判断通信成本。

\begin{knowledgebox}{讲义补充：源 Slide 14 是并行策略地图}
这页把大模型并行分成三层：data parallel 解决吞吐和数据切分；ZeRO/FSDP 解决 data parallel 下状态重复；model parallel 再把模型内部的 layers、width、sequence、experts、context 切开。真实训练通常不是三选一，而是把这些策略组合成 3D/4D parallelism。
\end{knowledgebox}

